\phantomsection\label{fig:vol04:jin:northern-rule-southern-frontier}
\begin{center}
\begin{tikzpicture}[x=.88cm,y=.88cm,font=\small]
  \fill[paper] (0,0) rectangle (11.8,6.4);

  \fill[source!19] (.55,3.7) -- (5.15,3.5) -- (5.25,5.85) -- (.75,5.9) -- cycle;
  \node[align=center] at (2.75,4.95) {女真兴起地带\\部族、河流与东北城镇};
  \fill[timeline!19] (5.6,3.7) -- (11.15,3.85) -- (11.2,5.85) -- (5.8,5.75) -- cycle;
  \node[align=center] at (8.65,4.95) {燕云、华北与中都\\1125年后重组的统治空间};
  \fill[dynasty!18] (.55,.65) -- (11.15,.75) -- (11.2,3.3) -- (5.55,3.4) -- (.65,3.35) -- cycle;
  \node[align=center] at (7.65,1.22) {河南、淮河与江南\\宋金边境、迁徙和财政压力};

  \node[draw=source,fill=paper,rounded corners=2pt,align=center,inner sep=3pt] (aguda)
    at (2.35,4.45) {1115\\阿骨打建国};
  \node[draw=timeline,fill=paper,rounded corners=2pt,align=center,inner sep=3pt] (yanjing)
    at (6.1,4.55) {燕京／中都\\道路、官署与都城};
  \node[draw=dynasty,fill=paper,rounded corners=2pt,align=center,inner sep=3pt] (kaifeng)
    at (8.7,2.55) {开封\\1214年后南迁中枢};
  \node[draw=dynasty,fill=paper,rounded corners=2pt,align=center,inner sep=3pt] (huai)
    at (6.05,2.1) {淮河--江淮\\宋金接触与战区};
  \node[draw=timeline,fill=paper,rounded corners=2pt,align=center,inner sep=3pt] (song)
    at (8.85,1.15) {南宋行在与江南\\财政、漕运网络};
  \node[draw=source,fill=paper,rounded corners=2pt,align=center,inner sep=3pt] (mongol)
    at (10.1,4.65) {1211--1234\\蒙古战争与金末};

  \draw[->,ultra thick,color=source!85] (aguda) -- node[above,sloped,font=\scriptsize]
    {灭辽、入燕的转换} (yanjing);
  \draw[->,ultra thick,color=dynasty!85] (yanjing) -- node[right,font=\scriptsize]
    {都城南迁与防务压力} (kaifeng);
  \draw[<->,thick,dashed,color=timeline!85] (kaifeng) -- node[right,font=\scriptsize]
    {和议、贸易与战争} (song);
  \draw[->,ultra thick,color=source!85] (mongol) -- node[above,sloped,font=\scriptsize]
    {围城、道路与降附者} (kaifeng);
  \draw[->,thick,dashed,color=timeline!85] (huai) -- (song);

  \node[draw=source!70,fill=paper,rounded corners=2pt,align=center,inner sep=3pt] at (2.55,2.35)
    {迁徙、户籍、军户与地方社会\\不能从政权转换一笔推得};
  \node[text=source,align=center,font=\scriptsize] at (6.25,.72)
    {范围、边界、箭头和城址位置均为约略示意；不表示精确辖境、连续战线、人口规模或族群分布。};
\end{tikzpicture}

\smallskip
\small\textit{图\quad 金的东北起源、华北统治与南方边境（1115--1234，示意）：据《金史》〈太祖本纪〉、〈世宗本纪〉、〈哀宗本纪〉及地理、兵、食货诸志，与《宋史》本纪、《三朝北盟会编》《建炎以来系年要录》所见宋金交涉互读；并参照女真、汉文、契丹文碑志、都城与边城考古及地方文书研究。图把建国、燕云转换、中都、开封南迁、淮河接触和蒙古战争并置，说明统治空间依赖道路、城镇、赋役与迁徙的重组；它不重绘任何时点的精确疆界、兵力、人口或地方服从。}
\end{center}
